\section{Related work}
\label{sec:related}

Four lines of work converge on the question this paper asks: whether a
measurable spectral correction translates into a measurable improvement.

\subsection{Spectral characterisation of generative models}

Generated images differ from real ones in their Fourier statistics, and the
difference is systematic. \citet{durall2020watch} traced a high-frequency
deficit in GAN samples to the up-convolution operators in the decoder;
\citet{dzanic2020fourier} showed the discrepancy is consistent enough across
architectures to serve as a detection signal, and \citet{chandrasegaran2021closer}
examined how much of it is intrinsic to the model rather than an artefact of
the last upsampling layer. \citet{wang2025spectral} supply the account of why
the deficit exists in the first place: solving the gradient-flow dynamics of a
denoiser yields an inverse-variance spectral law in which a Fourier mode's
learning time scales as $\tau \propto \lambda^{-1}$, so low-variance fine
detail is acquired orders of magnitude later than coarse structure. A
high-frequency deficit is therefore what an under-converged diffusion model
should look like, which is a competing explanation any frequency-aware
objective has to displace rather than assume. \citet{schwarz2021frequency} isolated generator and
discriminator separately and found the bias is not one effect: upsampling
operators bias the generator's spectrum, but checkerboard artefacts alone do
not explain the discrepancy, and the discriminator's difficulty is with
frequencies of \emph{low magnitude} rather than with high frequencies as such.

This literature establishes the phenomenon and, largely implicitly, a
motivation: because the spectrum is wrong, correcting it should help. That
implication is rarely tested directly, and it is what \S\ref{sec:worth} puts a
number on.

\subsection{Frequency-aware objectives and architectures}

Two families respond to the observation. The first modifies the
\emph{objective}. Focal Frequency Loss \citep{jiang2021focal} weights Fourier
components by how badly they are currently reconstructed, dynamically
emphasising the hard ones. Wavelet variants replace the Fourier basis with a
multi-resolution one, gaining spatial localisation: WaveDM \citep{huang2024wavedm}
applies wavelet supervision to restoration, and the objective audited here
weights wavelet sub-bands on a schedule tied to $\sigma_t$.

The second modifies the \emph{architecture} or the sampling path. WaveDiff
\citep{phung2023wavelet} moves the diffusion process itself into the wavelet
domain, primarily for speed; FreeU \citep{si2024freeu} rebalances the
backbone and skip contributions inside the U-Net at inference, and
\citet{yu2025frequency} rescale wavelet sub-bands during sampling to mitigate
exposure bias, and \citet{oh2026attention} edit cross-attention logits in the
Fourier domain at inference, one level deeper again. All three change the
frequency content of samples without touching training, which makes them the
closest relatives of the measurement in \S\ref{sec:worth} --- though there the
operation is used to price a mechanism rather than to improve one.
\citet{jiralerspong2025shaping} take a third route, shaping the frequency
content through the noising operator itself.

A recent comparative study of diffusion objectives \citep{kumar2025loss}
unifies several losses under the variational bound; frequency-weighted
objectives sit outside that unification, since their weights are not derived
from the bound but chosen to target a statistic of the output.

The family remains active. \citet{chandran2026spectral} add differentiable
Fourier \emph{and} wavelet losses to diffusion training as soft inductive
biases for frequency balance --- the closest contemporary sibling of the
objective we audit --- and report their largest gains on higher-resolution
unconditional data, where fine structure is hardest to model. That is worth
noting twice: it is independent support for the resolution caveat we state in
\S\ref{sec:limitations}, and it is a prediction the test of \S\ref{sec:worth}
could settle directly.

Reported gains in this family are typically small in absolute terms. The
improvements in \citet{chandran2026spectral} range from \num{0.02} to
\num{0.07} FID against baselines near \num{1.8}--\num{2.5}, averaged over three
seeds; those in the objective we audit are of comparable relative size. Our
results suggest that at that scale the question of whether an effect is present
at all requires the seed variance to be quantified first, and that even a
genuine spectral correction can be outweighed by what obtaining it costs.

\subsection{Timestep reweighting}

A parallel line reweights the loss across the noise schedule rather than across
frequency. \citet{nichol2021improved} resample timesteps by loss magnitude;
P2 \citep{choi2022perception} up-weights the region of the schedule where
perceptually relevant content is formed; Min-SNR \citep{hang2023efficient}
treats the schedule as a multi-task problem and clamps the per-timestep weight
by signal-to-noise ratio; \citet{karras2022elucidating} recast the whole design
space, weighting included, and \citet{kingma2021variational} derive a weighting
from the variational bound. \citet{li2026smoothed} unify the family from a
divergence-modification view and give an explicit cross-walk to Min-SNR, which
is what makes it meaningful to say a schedule \emph{is} a member of it rather
than merely correlated with one.

These two axes --- frequency and timestep --- are usually presented as
orthogonal. \S\ref{sec:confound} shows they are not orthogonal in the
formulation we audit: its frequency parameter simultaneously imposes a timestep
schedule correlated with Min-SNR at $|r| = 0.82$, so a result attributed to
frequency balance may belong to the other axis entirely. The multi-task framing
also supplies the natural alternative to hand-designed weights --- uncertainty
weighting \citep{kendall2018multi} and GradNorm \citep{chen2018gradnorm} --- which
\S\ref{sec:generality} evaluates.

\subsection{What the metrics can see}

A separate literature questions whether these metrics register what we assume.
\citet{barratt2018note} catalogued the Inception Score's failure modes;
\citet{chong2020effectively} showed FID is biased at finite sample size in a way
that depends on the model being evaluated; \citet{kynkaanniemi2023role}
demonstrated that FID responds strongly to the ImageNet class distribution of
the generated set, a property unrelated to image quality; and
\citet{stein2023exposing} found that the standard metrics rank diffusion models
against GANs in ways human raters do not reproduce.

Our contribution to this line is quantitative and specific: we measure how much
FID moves for a known spectral perturbation, and find that the answer depends
by more than an order of magnitude on whether the calibration is performed at
the metric's minimum or at a model's actual operating point
(\S\ref{sec:worth}).

\subsection{Replication and variance}

Finally, the practice this paper follows. \citet{bouthillier2021accounting}
showed that benchmark conclusions in machine learning frequently do not survive
proper accounting for sources of variance, seeds among them, and recommend
reporting comparisons with the variance made explicit.

\citet{dufour2026fid} put a number on it for exactly our setting. Across
several hundred networks trained on class-conditional ImageNet at
$256\times256$, retraining under an identical recipe with a different seed
moves FID $3.2\times$ more than redrawing samples from a fixed network, and
the coefficient of variation stays inside a \SIrange{1}{2}{\percent} band
regardless of model size or compute. They recommend treating any FID gap below
roughly \SI{1.3}{\percent} as inconclusive and reporting an error bar over
training seeds.

The agreement with our own measurement is worth stating because nothing about
it was arranged: our five-seed standard deviation of \num{0.21} at a baseline
FID of \fidMSE{} is a coefficient of variation of \SI{1.14}{\percent}, inside
their band, obtained on a different architecture, dataset and resolution. Two
independent estimates of how much a training seed is worth agreeing to within
a few tenths of a percent is the strongest evidence available that the quantity
is a property of diffusion training rather than of either setup. Differences of
the size routinely reported in the frequency-aware literature sit at or below
it, which is the methodological premise of \S\ref{sec:protocol}.

The 2026 sweep found no work that prices a training objective's stated
mechanism against a post-processed baseline, which remains this paper's
methodological contribution. Every cited entry has been verified against the
publisher or arXiv record.
